\documentclass[french]{book}
\usepackage{teach}
\usepackage{verbatim}
\usepackage{amsmath,amsfonts,amssymb}
\usepackage{pgfplots}
\RequirePackage{graphicx}
\RequirePackage{ifpdf}
 \ifpdf
   \DeclareGraphicsRule{*}{mps}{*}{}
 \fi
\usepackage{fancyhdr}
\usepackage{pstricks}
\usepackage{amsmath}
\usepackage{amssymb}
\usepackage{amsfonts}
\usepackage{amsthm}
\usepackage{mathrsfs}
\usepackage{array}
\usepackage{multicol}
\setlength{\multicolsep}{3pt}
\usepackage{graphicx}
\graphicspath{{Figures/}}
\usepackage{pstricks,pst-plot,pst-text,pst-tree,pst-eucl}
\usepackage[all]{xy}
\usepackage{fancybox}
\usepackage{color}
\usepackage{hyperref}
%  \usepackage[frenchb]{babel}
\usepackage{subfig}
\pgfplotsset{compat=newest}
%\usepackage{geometry}
%\geometry{top=1cm, bottom=1cm, left=2cm, right=2cm}
\usepackage[utf8]{inputenc}
\usepackage[french]{babel}
\redefinecolor{thm}{title}{bgcolor}{alizarin}
\redefinecolor{prop}{title}{bgcolor}{meatbrown}
\redefinecolor{defn}{title}{bgcolor}{olivine}
\definecolor{alizarin}{rgb}{0.82, 0.1, 0.26}
\definecolor{meatbrown}{rgb}{0.9, 0.72, 0.23}
\definecolor{olivine}{rgb}{0.6, 0.73, 0.45}
\usepackage{mathrsfs}
%%
\usepackage{yhmath} 
\newcommand{\R}{\mathbb {R}}
%\newcommand{\C}{\mathds {C}}
\newcommand{\N}{\mathbb {N}}
\newcommand{\Z}{\mathbb {Z}}
\newcommand{\Q}{\mathbb {Q}}
\newcommand{\D}{\mathbb {D}}
%%
\newcommand{\se}{\geq}
\newcommand{\ie}{\leq}
\newcommand{\tm}{\times}

%\newcommand{\ell}{l}
%\newcommand{\ell'}{l'}
%%
\newcommand{\ds}{\displaystyle}
\newcommand{\limn}{\lim_{n\rightarrow +\infty}}
\newcommand{\limo}[1][x]{\ds\lim_{#1\rightarrow 0}}
\newcommand{\limite}[2][x]{\ds\lim_{#1\rightarrow #2}}
\newcommand{\limpinf}[1][x]{\ds\lim_{#1\rightarrow +\infty}}
\newcommand{\limminf}[1][x]{\ds\lim_{#1\rightarrow -\infty}}
\newcommand{\limiteg}[2][x]{\ds\lim_{#1\xrightarrow{<}#2}}
\newcommand{\limited}[2][x]{\ds\lim_{#1\xrightarrow{>}#2}}
%%
\newcommand{\vect}[1]{\overrightarrow{#1}}
\newcommand{\arc}[1]{\wideparen{#1}}
\newcommand{\oij}{(\textit{O},{\overrightarrow{i}},{\overrightarrow{j}})}
%
\newcommand{\cp}[2]{%
        \begin{pmatrix}
        #1\\
        #2
        \end{pmatrix}%
        }
%
\newcommand{\calb}{\mathscr{B}}
\newcommand{\calc}{\mathscr{C}}
\newcommand{\cald}{\mathscr{D}}
\newcommand{\cale}{\mathscr{E}}
\newcommand{\calf}{\mathscr{F}}
\newcommand{\calp}{\mathscr{P}}
\newcommand{\cals}{\mathscr{S}}
\newcommand{\calt}{\mathscr{T}}
%
\newcommand{\suite}[1][u]{\left( #1_{n} \right)_{n\in\N}}
\newcommand{\suitep}[1][u]{\left( #1_{n} \right)_{n\in\N^{*}}}
\usetikzlibrary{arrows}

\newcommand{\poly}[1][a]{#1_0+#1_1x+\cdots+#1_nx^n}


\begin{document}
\tableofcontents
\chapter{Angles orientés}
\section{Définitions}
Une unité de longueur est choisie.
\subsection{Cercle trigonométrique, mesures d'un arc}
\begin{minipage}{10cm}
\begin{defn}
On appelle cercle trigonométrique un cercle de rayon 1, muni d'une origine et orienté dans le sens inverse des aiguilles d'une montre. Ce sens est appelé sens direct ou sens trigonométrique.
\par Soit $\calc$ un cercle trigonométrique de centre $O$ et d'origine $A$. À tout réel $x$, on peut associer un unique point $M$ du cercle en \og enroulant\fg{} la droite des réels autour du cercle. On dit que $x$ est une mesure de l'arc orienté $\arc{AM}$.
\end{defn}

\begin{rem}
Un point correspond à une infinité de réels.
\end{rem}

\begin{exemple}
Le point $A$ est associé aux réels 0, $2\pi$, $4\pi$, $-2\pi$...
Le point $B$ est associé aux réels $\dfrac{\pi}{2}$, $\dfrac{5\pi}{2}$, $-\dfrac{3\pi}{2}$...
Le point $A'$ est associé aux réels $\pi$, $3\pi$, $-\pi$...
\end{exemple}
\end{minipage}
\hfill
\begin{minipage}{4.6cm}
\psset{unit=1.3cm,PointSymbol=none}
\begin{center}
\begin{pspicture}(-2,-3)(1.5,3)
\pstGeonode[PosAngle=-135]{O}
\pstGeonode[PosAngle={0,90,180,-90},PointNameSep=0.5em](1,0){A}(0,1){B}(-1,0){A'}(0,-1){B'}
\pscircle(0,0)1
\pstLineAB{A'}{A}
\pstLineAB{B'}{B}
\pstGeonode[PointName=none](0.707,0.707){a}(-1.2,2.614){b}(-0.707,0.707){c}(-2,-0.586){d}
\pstOrtSym[PointName=none]{A'}{A}{a,b,c,d}
\psline{->}(1,-3)(1,3)
%\psline[linestyle=dashed]{->}(1,-3)(1,3)
\pstLineAB[linestyle=dashed]{->}{a}{b}
\pstLineAB[linewidth=1.8\pslinewidth]{->}{c}{d}
\pstLineAB[linestyle=dashed]{a'}{b'}
\pstLineAB[linewidth=1.8\pslinewidth]{c'}{d'}
\pstArcOAB[linewidth=1.8\pslinewidth]{O}{c'}{c}
\pstGeonode[PointSymbol=|,dotangle=90,dotscale=1.3](1,2.094){x}
\pstGeonode[PointSymbol=|,dotangle=90,dotscale=1.3,PointName=1](1,1){i}
\pstGeonode[PointSymbol=|,dotangle=-45,dotscale=1.3,PointName=x,PosAngle=15](-0.218,1.633){x'}
\pstGeonode[PointSymbol=|,dotangle=30,dotscale=1.3,PosAngle=150](-0.5,0.866){M}
\ncarc[nodesep=3mm,arcangle=-30]{->}{x}{x'}
\ncarc[nodesep=2mm,arcangle=-30]{->}{x'}{M}
\end{pspicture}
\end{center}
\end{minipage}

\subsection{Mesures d'un angle orienté de vecteurs}
Étant donnés deux vecteurs non nuls $\vect{u}$ et $\vect{v}$ et un cercle trigonométrique $\calc$ de centre $O$, on appelle $A$ le point d'intersection de $\calc$ et de la droite $(O,\vect{u})$ et $B$ le point d'intersection de $\calc$ et de la droite $(O,\vect{v})$.

\begin{minipage}{10.5cm}
\begin{defn}
On appelle mesure de l'angle orienté $(\vect{u},\vect{v})$ toute mesure de l'arc orienté $\arc{AB}$.

Si $x$ est une de ces mesures, toute autre mesure s'écrit $y=x+2k\pi$ avec $k\in\Z$.

On note (de façon abusive) : $(\vect{u},\vect{v})=x+k\times 2\pi, k\in\Z$.
\end{defn}
\end{minipage}
\hfill
\begin{minipage}{4cm}
	\begin{center}
	\includegraphics{CoursAnglesFigures.1}
	\end{center}
\end{minipage}

% Remarque : Si $x$ est une mesure de $(\vect{u},\vect{v})$, toute autre mesure s'écrit $y=x+2k\pi$ avec $k\in\Z$. On note donc aussi $(\vect{u},\vect{v})=x+2k\pi, \; k\in\Z$.

\begin{defn}
On appelle mesure principale de $(\vect{u},\vect{v})$, l'unique mesure appartenant à $]-\pi ; \pi]$.

\vspace{4mm}
\end{defn}

\begin{exemple}
Déterminer la mesure principale d'un angle dont une mesure est $\dfrac{126\pi}{5}$
\end{exemple}


\[\dfrac{126\pi}{5} = \dfrac{12\times10\pi+6\pi}{5} = 12\times2\pi+\dfrac{6\pi}{5} = 12\times2\pi+2\pi-\dfrac{4\pi}{5} = -\dfrac{4\pi}{5}+13\times 2\pi\]
La mesure principale de cet angle est donc $\dfrac{-4\pi}{5}$.


\section{Propriétés}
\subsection{Angles et colinéarité}
\begin{prop}
Soient $\vect{u}$ et $\vect{v}$ deux vecteurs non nuls.
\par $\vect{u}$ et $\vect{v}$ sont colinéaires et de même sens si et seulement si $(\vect{u},\vect{v})=0+2k\pi$.
\par $\vect{u}$ et $\vect{v}$ sont colinéaires et de sens contraires si et seulement si $(\vect{u},\vect{v})=\pi+2k\pi$.
\end{prop}

\begin{center}
\includegraphics{CoursAnglesFigures.2}
\hspace{3em}
\includegraphics{CoursAnglesFigures.3}
\end{center}


\subsection{Relation de Chasles}
\begin{minipage}{10cm}
\begin{prop}
Pour tous vecteurs non nuls $\vect{u}$, $\vect{v}$ et $\vect{w}$ : \[(\vect{u},\vect{v})+(\vect{v},\vect{w})=(\vect{u},\vect{w})+2k\pi\]
\end{prop}
\end{minipage}
\hfill
\begin{minipage}{4cm}
\includegraphics{CoursAnglesFigures.4}
\end{minipage}


\begin{exemple}
Dans la figure suivante, démontrer que les droites $(AB)$ et $(CD)$ sont parallèles.
\end{exemple}

\begin{minipage}{9.5cm}

\begin{multline*}
(\vect{AB},\vect{CD})=(\vect{AB},\vect{BC})+(\vect{BC},\vect{CD})+2k\pi\\
=\dfrac{\pi}{3}+\dfrac{2\pi}{3}+2k\pi=\pi+2k\pi
\end{multline*}
\par Les vecteurs $\vect{AB}$ et $\vect{CD}$ sont donc colinéaires.
\par \underline{Conclusion} : $(AB)$ et $(CD)$ sont parallèles.

\end{minipage}
\hfill
\begin{minipage}{5cm}
\begin{center}
\includegraphics{CoursAnglesFigures.9}
\end{center}
\end{minipage}

\newpage

\begin{prop}
\par Pour tous vecteurs non nuls $\vect{u}$ et $\vect{v}$ :

\vspace{1ex}
\begin{tabular}{ll@{\hspace{5em}}ll}
$(\vect{v},\vect{u})=-(\vect{u},\vect{v})+2k\pi$ & $(1)$ & $(\vect{u},-\vect{v})=(\vect{u},\vect{v})+\pi+2k\pi$ & $(2)$ \\[1ex]
$(-\vect{u},\vect{v})=(\vect{u},\vect{v})+\pi+2k\pi$ & $(3)$ & $(-\vect{u},-\vect{v})=(\vect{u},\vect{v})+2k\pi$ & $(4)$
\end{tabular}
\end{prop}

\includegraphics{CoursAnglesFigures.5}
\hfill
\includegraphics{CoursAnglesFigures.6}
\hfill
\includegraphics{CoursAnglesFigures.7}
\hfill
\includegraphics{CoursAnglesFigures.8}

\begin{demo}
\par (1) : $(\vect{v},\vect{u})+(\vect{u},\vect{v})=(\vect{v},\vect{v})+2k\pi=2k\pi$ ainsi $(\vect{v},\vect{u})=-(\vect{u},\vect{v})+2k\pi$
\par (2) : $(\vect{u},-\vect{v})=(\vect{u},\vect{v})+(\vect{v},-\vect{v})+2k\pi=(\vect{u},\vect{v})+\pi+2k\pi$
\par (3) se démontre comme (2).

\begin{tabbing}
(4) : $(-\vect{u},-\vect{v})$ \= $=(-\vect{u},\vect{v})+\pi+2k\pi$ d'après (2)\\
 \> $= (\vect{u},\vect{v})+\pi+\pi+2k\pi$ d'après (3)\\
ainsi $(-\vect{u},-\vect{v})=(\vect{u},\vect{v})+2k\pi$
\end{tabbing}
\vspace{-1em}
\end{demo}

\end{document}
